\chapter{Continuous Dynamical Systems}\label{ch:3}

In Chapter 2, we studied discrete dynamical systems generated by iterating a map $T: X \to X$. We now turn to the continuous-time setting, where the state of a system evolves according to an ordinary differential equation (ODE). This is the natural framework for modeling physical systems governed by Newton's laws, population dynamics in ecology, chemical reaction kinetics, and electrical circuits, among many others.

The shift from discrete to continuous time is more than a change of notation. New phenomena arise---limit cycles, for instance, have no direct analogue in one-dimensional discrete systems---and certain phenomena disappear: we shall see that chaos is impossible for flows on the plane, a striking contrast with the rich dynamics of one-dimensional maps.


\bigskip


\section{ODEs as Dynamical Systems}\label{sec:3.1}

\subsection{Setup and Terminology}\label{sec:3.1.1}

An \textbf{autonomous ordinary differential equation} on $\mathbb{R}^n$ is a system of the form

\begin{equation}\label{eq:3.1}\tag{3.1}
\dot{x} = f(x),
\end{equation}

where $x \in \mathbb{R}^n$, $\dot{x} = dx/dt$, and $f: U \to \mathbb{R}^n$ is a vector field defined on some open set $U \subseteq \mathbb{R}^n$. The word \textit{autonomous} means that $f$ does not depend explicitly on time; the law governing the evolution depends only on the current state.

A \textbf{solution} (or \textbf{trajectory}, or \textbf{integral curve}) of (3.1) with initial condition $x(0) = x_0$ is a differentiable function $x: I \to \mathbb{R}^n$, defined on some interval $I$ containing $0$, such that $x'(t) = f(x(t))$ for all $t \in I$ and $x(0) = x_0$.

A point $x^* \in U$ is called a \textbf{fixed point} (or \textbf{equilibrium}, or \textbf{critical point}) of the system if $f(x^*) = 0$. At such a point, the constant function $x(t) = x^*$ for all $t$ is a solution.

\subsection{The Flow}\label{sec:3.1.2}

When (3.1) has a unique solution for each initial condition $x_0$, we can define the \textbf{flow} of the system. Write $\phi_t(x_0)$ for the value at time $t$ of the solution starting at $x_0$. This gives a family of maps

\[
\phi_t: U \to \mathbb{R}^n, \quad t \in \mathbb{R},
\]

satisfying the following properties:

\begin{enumerate}[nosep]
  \item \textbf{Identity:} $\phi_0(x) = x$ for all $x \in U$.
  \item \textbf{Group property:} $\phi_{t+s}(x) = \phi_t(\phi_s(x))$ for all $s, t$ for which both sides are defined.
  \item \textbf{Smoothness:} If $f$ is $C^k$, then $\phi_t(x)$ is $C^k$ in both $t$ and $x$.
\end{enumerate}

The group property deserves emphasis. It says that flowing for time $s$ and then for time $t$ is the same as flowing for time $t + s$. In the language of algebra, $\{\phi_t\}_{t \in \mathbb{R}}$ is a one-parameter group of transformations of $U$. This is precisely the continuous analogue of iterating a map: in a discrete system, $T^{n+m} = T^n \circ T^m$.

\begin{proof}[Proof of the group property]
of the group property.** Fix $x_0 \in U$ and $s \in \mathbb{R}$. Define two functions of $t$:

\begin{itemize}[nosep]
  \item $\alpha(t) = \phi_{t+s}(x_0)$,
  \item $\beta(t) = \phi_t(\phi_s(x_0))$.
\end{itemize}

Both satisfy the ODE $\dot{y} = f(y)$ (for $\alpha$, differentiate with respect to $t$ and use the definition of flow; for $\beta$, this holds by definition). At $t = 0$, both equal $\phi_s(x_0)$. By uniqueness of solutions, $\alpha(t) = \beta(t)$ for all $t$.
\end{proof}


Note the essential role of uniqueness here. Without it, the flow is not well-defined, and the entire framework breaks down.

\subsection{Continuous vs. Discrete}\label{sec:3.1.3}

It is worth pausing to compare the two settings.

| | Discrete: $x_{n+1} = T(x_n)$ | Continuous: $\dot{x} = f(x)$ |
|---|---|---|
| Time | $\mathbb{Z}$ or $\mathbb{N}$ | $\mathbb{R}$ |
| Evolution operator | $T^n$ (iterate) | $\phi_t$ (flow) |
| Group property | $T^{n+m} = T^n \circ T^m$ | $\phi_{t+s} = \phi_t \circ \phi_s$ |
| Orbits | Sequences $\{x_n\}$ | Curves $\{x(t)\}$ |

One can always pass from continuous to discrete by \textbf{sampling} the flow: set $T = \phi_\tau$ for some fixed time step $\tau > 0$. The map $T$ is then called the **time-$\tau$ map** of the flow. This construction will reappear when we discuss Poincar\'e sections and, later, in the context of reservoir computing.


\bigskip


\section{Existence and Uniqueness}\label{sec:3.2}

\subsection{The Picard--Lindel\"of Theorem}\label{sec:3.2.1}

The fundamental theorem guaranteeing that solutions exist and are unique is the following.

\textbf{Theorem 3.1 (Picard--Lindel\"of).} *Let $U \subseteq \mathbb{R}^n$ be open, and let $f: U \to \mathbb{R}^n$ be continuous. Suppose $f$ is \textbf{locally Lipschitz} on $U$: for every compact set $K \subset U$, there exists a constant $L > 0$ such that*

\[
\|f(x) - f(y)\| \leq L\|x - y\| \quad \text{for all } x, y \in K.
\]

*Then for every $x_0 \in U$, there exists an interval $(-\varepsilon, \varepsilon)$ with $\varepsilon > 0$ and a unique solution $x: (-\varepsilon, \varepsilon) \to U$ of $\dot{x} = f(x)$ with $x(0) = x_0$.*

Moreover, this solution can be extended to a \textbf{maximal interval of existence} $(t_{\min}, t_{\max})$ with $-\infty \leq t_{\min} < 0 < t_{\max} \leq +\infty$, and if $t_{\max} < \infty$, then the solution eventually leaves every compact subset of $U$ as $t \to t_{\max}^-$ (and similarly for $t_{\min}$).

\textbf{Sufficient condition.} If $f$ is $C^1$ (i.e., continuously differentiable), then $f$ is locally Lipschitz. This follows immediately from the mean value theorem. In practice, the vector fields one encounters in applications are almost always $C^1$ or better, so the Picard--Lindel\"of theorem applies.

\begin{proof}[Proof sketch]
sketch.** The standard proof proceeds by converting the ODE to an integral equation

\[
x(t) = x_0 + \int_0^t f(x(s))\, ds
\]

and showing that the map $\mathcal{T}$ defined by $(\mathcal{T} x)(t) = x_0 + \int_0^t f(x(s))\,ds$ is a contraction on an appropriate function space (continuous functions on $[-\varepsilon, \varepsilon]$ with values in a ball around $x_0$, equipped with the supremum norm). The Banach fixed-point theorem then gives existence and uniqueness of a fixed point, which is the desired solution.

In detail: choose $r > 0$ so that the closed ball $\overline{B}_r(x_0) \subset U$, and let $M = \sup_{\overline{B}_r(x_0)} \|f\|$, $L$ the Lipschitz constant on $\overline{B}_r(x_0)$. Set $\varepsilon = \min(r/M, 1/(2L))$. On the space $C([-\varepsilon,\varepsilon], \overline{B}_r(x_0))$ with the supremum norm, one checks:

\begin{enumerate}[nosep]
  \item $\mathcal{T}$ maps this space to itself (because $\|\mathcal{T}x(t) - x_0\| \leq M\varepsilon \leq r$).
  \item $\|\mathcal{T}x - \mathcal{T}y\|_\infty \leq L\varepsilon \|x - y\|_\infty \leq \frac{1}{2}\|x - y\|_\infty$.
\end{enumerate}

By the Banach fixed-point theorem, $\mathcal{T}$ has a unique fixed point.
\end{proof}


\subsection{Failure of Uniqueness}\label{sec:3.2.2}

What happens when $f$ is continuous but not Lipschitz? Consider the scalar ODE

\begin{equation}\label{eq:3.2}\tag{3.2}
\dot{x} = x^{2/3}, \quad x(0) = 0.
\end{equation}

The vector field $f(x) = x^{2/3}$ is continuous on $\mathbb{R}$ but is \textit{not} Lipschitz at $x = 0$: for $x > 0$, $f'(x) = \frac{2}{3}x^{-1/3} \to \infty$ as $x \to 0^+$.

One can verify by direct substitution that for any $c \geq 0$, the function

\[
x_c(t) = \begin{cases} 0 & \text{if } t \leq c, \\ \left(\frac{t - c}{3}\right)^3 & \text{if } t > c, \end{cases}
\]

is a solution of (3.2) with $x_c(0) = 0$. (Check: for $t > c$, $\dot{x}_c = 3 \cdot \frac{1}{27}(t-c)^2 = \frac{(t-c)^2}{9}$, while $x_c^{2/3} = \left(\frac{(t-c)^3}{27}\right)^{2/3} = \frac{(t-c)^2}{9}$.) The trivial solution $x(t) \equiv 0$ is also valid. So uniqueness fails spectacularly: there are uncountably many solutions through the origin.

\subsection{The Peano Existence Theorem}\label{sec:3.2.3}

Even without the Lipschitz condition, continuity alone guarantees existence (though not uniqueness).

\textbf{Theorem 3.2 (Peano).} *Let $f: U \to \mathbb{R}^n$ be continuous on an open set $U \subseteq \mathbb{R}^n$. Then for every $x_0 \in U$, there exists $\varepsilon > 0$ and a solution $x: [0, \varepsilon) \to U$ of $\dot{x} = f(x)$ with $x(0) = x_0$.*

The proof uses the Arzel\`a--Ascoli theorem and is more delicate than the Picard--Lindel\"of proof. Since we will always assume at least Lipschitz regularity in what follows, we state Peano's theorem without proof and refer the reader to Coddington and Levinson (1955) for details.


\bigskip


\section{Linear Systems}\label{sec:3.3}

\subsection{The Matrix Exponential}\label{sec:3.3.1}

The simplest class of autonomous ODEs is the \textbf{linear} case:

\begin{equation}\label{eq:3.3}\tag{3.3}
\dot{x} = Ax,
\end{equation}

where $A$ is a real $n \times n$ matrix and $x \in \mathbb{R}^n$. These systems can be solved exactly.

\begin{definition}\label{def:3.3.}
The \textbf{matrix exponential} of $A \in \mathbb{R}^{n \times n}$ is

\[
e^A = \sum_{k=0}^{\infty} \frac{A^k}{k!} = I + A + \frac{A^2}{2!} + \frac{A^3}{3!} + \cdots
\]

This series converges absolutely for every $A$ (with respect to any matrix norm), since $\|A^k/k!\| \leq \|A\|^k/k!$ and $\sum \|A\|^k/k! = e^{\|A\|} < \infty$.
\end{definition}


\begin{proposition}\label{prop:3.4.}
*The solution of $\dot{x} = Ax$, $x(0) = x_0$ is*

\[
x(t) = e^{At} x_0.
\]
\end{proposition}


\begin{proof}
* Differentiate term by term:

\[
\frac{d}{dt} e^{At} = \frac{d}{dt} \sum_{k=0}^{\infty} \frac{(At)^k}{k!} = \sum_{k=1}^{\infty} \frac{k A^k t^{k-1}}{k!} = A \sum_{k=1}^{\infty} \frac{(At)^{k-1}}{(k-1)!} = A e^{At}.
\]

Term-by-term differentiation is justified by the uniform convergence of the series on compact intervals. So $x(t) = e^{At}x_0$ satisfies $\dot{x} = Ae^{At}x_0 = Ax(t)$, and $x(0) = e^0 x_0 = Ix_0 = x_0$.
\end{proof}


The flow of the linear system is therefore $\phi_t(x_0) = e^{At}x_0$. Key properties of the matrix exponential include:

\begin{itemize}[nosep]
  \item $e^{0} = I$.
  \item $e^{(t+s)A} = e^{tA} e^{sA}$ (confirming the group property).
  \item If $AB = BA$, then $e^{A+B} = e^A e^B$. \textbf{Caution:} this fails in general when $A$ and $B$ do not commute.
  \item $\det(e^A) = e^{\operatorname{tr}(A)}$ (Jacobi's formula). In particular, $e^A$ is always invertible.
  \item If $P$ is invertible, then $e^{PAP^{-1}} = P e^A P^{-1}$.
\end{itemize}

The last property is crucial: it reduces computation of $e^{At}$ to the case where $A$ is in a convenient canonical form.

\subsection{Computing the Matrix Exponential via Jordan Form}\label{sec:3.3.2}

Every real matrix $A$ can be brought (over $\mathbb{C}$) to \textbf{Jordan normal form}: there exists an invertible matrix $P$ such that $A = PJP^{-1}$, where $J$ is block-diagonal with Jordan blocks

\[
J_k(\lambda) = \begin{pmatrix} \lambda & 1 & 0 & \cdots & 0 \\ 0 & \lambda & 1 & \cdots & 0 \\ \vdots & & \ddots & \ddots & \vdots \\ 0 & 0 & \cdots & \lambda & 1 \\ 0 & 0 & \cdots & 0 & \lambda \end{pmatrix} \in \mathbb{C}^{k \times k}.
\]

Since $e^{PJP^{-1}t} = P e^{Jt} P^{-1}$ and $e^{Jt}$ is block-diagonal with blocks $e^{J_k(\lambda)t}$, we need only compute

\begin{equation}\label{eq:3.4}\tag{3.4}
e^{J_k(\lambda)t} = e^{\lambda t} \begin{pmatrix} 1 & t & t^2/2! & \cdots & t^{k-1}/(k-1)! \\ 0 & 1 & t & \cdots & t^{k-2}/(k-2)! \\ \vdots & & \ddots & \ddots & \vdots \\ 0 & 0 & \cdots & 1 & t \\ 0 & 0 & \cdots & 0 & 1 \end{pmatrix}.
\end{equation}

This follows by writing $J_k(\lambda) = \lambda I + N$ where $N$ is the nilpotent part (the superdiagonal of ones), noting that $\lambda I$ and $N$ commute, so $e^{J_k(\lambda)t} = e^{\lambda t} e^{Nt}$, and computing $e^{Nt}$ using the fact that $N^k = 0$.

For a diagonalizable matrix (when all Jordan blocks are $1 \times 1$), this simplifies: if $A = P \operatorname{diag}(\lambda_1, \ldots, \lambda_n) P^{-1}$, then $e^{At} = P \operatorname{diag}(e^{\lambda_1 t}, \ldots, e^{\lambda_n t}) P^{-1}$.

\subsection{Classification of Linear Flows in 2D}\label{sec:3.3.3}

Consider $\dot{x} = Ax$ with $A \in \mathbb{R}^{2 \times 2}$. The qualitative behavior of the flow is determined entirely by the eigenvalues of $A$. Let $\tau = \operatorname{tr}(A)$ and $\delta = \det(A)$. The eigenvalues satisfy $\lambda^2 - \tau\lambda + \delta = 0$, so $\lambda = \frac{1}{2}(\tau \pm \sqrt{\tau^2 - 4\delta})$. We classify:

**Case 1: Two distinct real eigenvalues $\lambda_1, \lambda_2$ ($\tau^2 - 4\delta > 0$).**

After a change of basis, the system becomes $\dot{y}_1 = \lambda_1 y_1$, $\dot{y}_2 = \lambda_2 y_2$, with solutions $y_1(t) = c_1 e^{\lambda_1 t}$, $y_2(t) = c_2 e^{\lambda_2 t}$. Trajectories satisfy $y_2 = C |y_1|^{\lambda_2/\lambda_1}$ for some constant $C$.

\begin{itemize}[nosep]
  \item \textbf{Stable node} ($\lambda_1 < \lambda_2 < 0$): All trajectories approach the origin as $t \to \infty$. They are tangent to the "slow" eigendirection (corresponding to the eigenvalue closer to zero, i.e., $\lambda_2$) as they approach. The phase portrait looks like a sink, with all arrows pointing inward.

  \item \textbf{Unstable node} ($0 < \lambda_1 < \lambda_2$): Time-reversed stable node. All trajectories escape to infinity. Tangent to the slow eigendirection near the origin.

  \item \textbf{Saddle point} ($\lambda_1 < 0 < \lambda_2$): The eigendirection for $\lambda_1$ is the \textbf{stable manifold} $W^s$---trajectories starting on it approach the origin. The eigendirection for $\lambda_2$ is the \textbf{unstable manifold} $W^u$---trajectories starting on it escape the origin. All other trajectories approach the unstable manifold as $t \to +\infty$ and the stable manifold as $t \to -\infty$. The phase portrait has a characteristic "X" shape.
\end{itemize}

**Case 2: Complex eigenvalues $\lambda = \alpha \pm i\beta$ with $\beta \neq 0$ ($\tau^2 - 4\delta < 0$).**

Here $\alpha = \tau/2$ and $\beta = \sqrt{4\delta - \tau^2}/2$. After a real change of basis, the system becomes

\[
\dot{y}_1 = \alpha y_1 - \beta y_2, \qquad \dot{y}_2 = \beta y_1 + \alpha y_2.
\]

Switching to polar coordinates $r = \sqrt{y_1^2 + y_2^2}$, $\theta = \arctan(y_2/y_1)$:

\[
\dot{r} = \alpha r, \qquad \dot{\theta} = \beta.
\]

So $r(t) = r_0 e^{\alpha t}$ and $\theta(t) = \theta_0 + \beta t$. Trajectories are logarithmic spirals (when $\alpha \neq 0$) or circles (when $\alpha = 0$).

\begin{itemize}[nosep]
  \item \textbf{Stable spiral (focus)} ($\alpha < 0$): Trajectories spiral inward toward the origin.
  \item \textbf{Unstable spiral (focus)} ($\alpha > 0$): Trajectories spiral outward.
  \item \textbf{Center} ($\alpha = 0$, i.e., $\tau = 0$): Trajectories are closed ellipses (circles in the eigenbasis). The origin is Lyapunov stable but not asymptotically stable.
\end{itemize}

**Case 3: Repeated eigenvalue $\lambda = \tau/2$ ($\tau^2 - 4\delta = 0$).**

\begin{itemize}[nosep]
  \item If $A$ is diagonalizable (i.e., $A = \lambda I$), then $e^{At} = e^{\lambda t}I$ and trajectories are rays from the origin. This is called a \textbf{star node}.
  \item If $A$ is not diagonalizable, the Jordan form is $J = \begin{pmatrix} \lambda & 1 \\ 0 & \lambda \end{pmatrix}$, giving $e^{Jt} = e^{\lambda t}\begin{pmatrix} 1 & t \\ 0 & 1 \end{pmatrix}$. Trajectories approach the origin tangent to the eigendirection. This is called a \textbf{degenerate (improper) node}.
\end{itemize}

**Case 4: Degenerate cases ($\delta = 0$).**

When $\det(A) = 0$, at least one eigenvalue is zero, and the origin is not an isolated fixed point. Instead, there is an entire line (or plane) of fixed points. We exclude this case from the classification since the origin is not an isolated equilibrium.

\textbf{Summary.} The classification is governed by the $(\tau, \delta)$-plane:

\begin{itemize}[nosep]
  \item The parabola $\tau^2 = 4\delta$ separates nodes from spirals/centers.
  \item The $\delta$-axis ($\tau = 0$, $\delta > 0$) corresponds to centers.
  \item The region $\delta < 0$ corresponds to saddle points.
  \item Stability is determined by the sign of $\tau$: stable if $\tau < 0$, unstable if $\tau > 0$.
\end{itemize}

\textbf{Worked Example 3.5.} Consider $A = \begin{pmatrix} -1 & 2 \\ 0 & -3 \end{pmatrix}$. The eigenvalues are $\lambda_1 = -1$ and $\lambda_2 = -3$ (read off the diagonal since $A$ is upper-triangular). Both are negative, so the origin is a \textbf{stable node}.

The eigenvectors: for $\lambda_1 = -1$, $(A + I)v = 0$ gives $\begin{pmatrix} 0 & 2 \\ 0 & -2 \end{pmatrix}v = 0$, so $v_1 = \begin{pmatrix}1\\0\end{pmatrix}$. For $\lambda_2 = -3$, $(A + 3I)v = 0$ gives $\begin{pmatrix}2 & 2 \\ 0 & 0\end{pmatrix}v = 0$, so $v_2 = \begin{pmatrix}1 \\ -1\end{pmatrix}$.

The general solution is $x(t) = c_1 e^{-t}\begin{pmatrix}1\\0\end{pmatrix} + c_2 e^{-3t}\begin{pmatrix}1\\-1\end{pmatrix}$. As $t \to \infty$, the $e^{-3t}$ term decays faster, so trajectories become tangent to the $v_1$ direction (the slow eigendirection).

\textbf{Worked Example 3.6.} Consider $A = \begin{pmatrix} 0 & -1 \\ 1 & 0 \end{pmatrix}$. The eigenvalues are $\lambda = \pm i$ (here $\tau = 0$, $\delta = 1$). The origin is a \textbf{center}. The solutions are $x_1(t) = c_1 \cos t - c_2 \sin t$, $x_2(t) = c_1 \sin t + c_2 \cos t$---uniform rotation. Trajectories are circles centered at the origin.


\bigskip


\section{Phase Portraits in 2D}\label{sec:3.4}

\subsection{Nullclines}\label{sec:3.4.1}

For a two-dimensional system

\[
\dot{x}_1 = f_1(x_1, x_2), \qquad \dot{x}_2 = f_2(x_1, x_2),
\]

the **$x_1$-nullcline** is the set $\{(x_1, x_2) : f_1(x_1, x_2) = 0\}$ and the **$x_2$-nullcline** is $\{(x_1, x_2) : f_2(x_1, x_2) = 0\}$.

On the $x_1$-nullcline, $\dot{x}_1 = 0$, so the flow is purely vertical. On the $x_2$-nullcline, $\dot{x}_2 = 0$, so the flow is purely horizontal. Fixed points lie at intersections of the two nullclines. Between the nullclines, the signs of $\dot{x}_1$ and $\dot{x}_2$ are constant, which determines the general direction of the flow in each region.

This gives a systematic method for sketching phase portraits:

\begin{enumerate}[nosep]
  \item Find the nullclines.
  \item Mark their intersections (fixed points).
  \item Determine the signs of $\dot{x}_1$ and $\dot{x}_2$ in each region.
  \item Indicate the flow direction on the nullclines (vertical on $x_1$-nullclines, horizontal on $x_2$-nullclines).
  \item Sketch representative trajectories consistent with this information.
\end{enumerate}

\subsection{Worked Example: The Simple Pendulum}\label{sec:3.4.2}

The equation of motion for a frictionless pendulum of unit length and unit mass under gravity $g = 1$ is

\begin{equation}\label{eq:3.5}\tag{3.5}
\ddot{\theta} + \sin\theta = 0.
\end{equation}

We convert this to a first-order system by setting $x_1 = \theta$ (angle) and $x_2 = \dot{\theta}$ (angular velocity):

\begin{equation}\label{eq:3.6}\tag{3.6}
\dot{x}_1 = x_2, \qquad \dot{x}_2 = -\sin x_1.
\end{equation}

\textbf{Fixed points.} Setting both right-hand sides to zero: $x_2 = 0$ and $\sin x_1 = 0$, so $x_1 = k\pi$ for $k \in \mathbb{Z}$. The physically distinct fixed points are $(0, 0)$ (pendulum hanging down) and $(\pi, 0)$ (pendulum balanced upright).

\textbf{Nullclines.} The $x_1$-nullcline is $x_2 = 0$ (the horizontal axis). The $x_2$-nullcline is $\sin x_1 = 0$, i.e., $x_1 = k\pi$ (vertical lines).

**Linearization at $(0, 0)$.** The Jacobian (see Section 3.5) is

\[
Df(0,0) = \begin{pmatrix} 0 & 1 \\ -\cos 0 & 0 \end{pmatrix} = \begin{pmatrix} 0 & 1 \\ -1 & 0 \end{pmatrix}.
\]

Eigenvalues: $\lambda = \pm i$. The origin is a \textbf{center} of the linearized system.

**Linearization at $(\pi, 0)$.** The Jacobian is

\[
Df(\pi, 0) = \begin{pmatrix} 0 & 1 \\ -\cos\pi & 0 \end{pmatrix} = \begin{pmatrix} 0 & 1 \\ 1 & 0 \end{pmatrix}.
\]

Eigenvalues: $\lambda = \pm 1$. This is a \textbf{saddle point}.

\textbf{Energy and the global phase portrait.} The system (3.6) is Hamiltonian with energy

\[
E(\theta, \dot\theta) = \frac{1}{2}\dot\theta^2 - \cos\theta = \frac{1}{2}x_2^2 - \cos x_1.
\]

One verifies that $\frac{d}{dt}E = x_2 \dot{x}_2 + \sin(x_1)\dot{x}_1 = x_2(-\sin x_1) + \sin(x_1)\cdot x_2 = 0$, so $E$ is conserved along trajectories. The level curves of $E$ are therefore exactly the trajectories.

At the saddle $(\pi, 0)$: $E(\pi, 0) = 0 + 1 = 1$. The level set $E = 1$ consists of the \textbf{separatrices} that connect the saddle points---these are called \textbf{heteroclinic orbits} (connecting $(-\pi, 0)$ to $(\pi, 0)$ and vice versa).

\begin{itemize}[nosep]
  \item For $E < 1$ near $(0,0)$: closed orbits (oscillations, the pendulum swings back and forth).
  \item For $E > 1$: open orbits (rotations, the pendulum goes over the top).
  \item $E = 1$: the separatrices.
\end{itemize}

This phase portrait is one of the canonical examples in dynamics. The coexistence of a center, saddle points, and the separatrix structure connecting them illustrates the richness that even simple mechanical systems exhibit.

\subsection{Worked Example: Lotka--Volterra Predator-Prey}\label{sec:3.4.3}

The classical predator-prey model is

\begin{equation}\label{eq:3.7}\tag{3.7}
\dot{x} = \alpha x - \beta x y, \qquad \dot{y} = \delta x y - \gamma y,
\end{equation}

where $x \geq 0$ is the prey population, $y \geq 0$ is the predator population, and $\alpha, \beta, \gamma, \delta > 0$ are parameters. For concreteness, take $\alpha = \beta = \gamma = \delta = 1$:

\begin{equation}\label{eq:3.8}\tag{3.8}
\dot{x} = x(1 - y), \qquad \dot{y} = y(x - 1).
\end{equation}

\textbf{Fixed points.} From $\dot{x} = 0$: $x = 0$ or $y = 1$. From $\dot{y} = 0$: $y = 0$ or $x = 1$. The fixed points (in the first quadrant and its boundary) are $(0, 0)$ and $(1, 1)$.

\textbf{Nullclines.} The $x$-nullcline consists of $x = 0$ (the $y$-axis) and $y = 1$ (a horizontal line). The $y$-nullcline consists of $y = 0$ (the $x$-axis) and $x = 1$ (a vertical line). These divide the first quadrant into four regions:

| Region | $\dot{x}$ | $\dot{y}$ | Flow direction |
|--------|-----------|-----------|----------------|
| $x < 1, y < 1$ | $+$ | $-$ | right and down |
| $x > 1, y < 1$ | $+$ | $+$ | right and up |
| $x > 1, y > 1$ | $-$ | $+$ | left and up |
| $x < 1, y > 1$ | $-$ | $-$ | left and down |

The flow circulates counterclockwise around $(1,1)$.

**Linearization at $(1,1)$.** The Jacobian of the right-hand side of (3.8) is

\[
J = \begin{pmatrix} 1 - y & -x \\ y & x - 1 \end{pmatrix}.
\]

At $(1,1)$:

\[
J(1,1) = \begin{pmatrix} 0 & -1 \\ 1 & 0 \end{pmatrix}.
\]

The eigenvalues are $\lambda = \pm i$, so the linearization predicts a \textbf{center}.

\textbf{Conserved quantity.} Like the pendulum, this system has a conserved quantity. Define

\[
H(x, y) = \delta x - \gamma \ln x + \beta y - \alpha \ln y.
\]

For our parameters: $H(x,y) = x - \ln x + y - \ln y$. One can verify $\dot{H} = 0$ along trajectories (exercise). Since $H$ is a strictly convex function with a unique minimum at $(1,1)$ (check: $H_x = 1 - 1/x = 0$ at $x=1$, $H_{xx} = 1/x^2 > 0$, similarly for $y$), its level sets are closed curves surrounding $(1,1)$. Therefore the trajectories in the first quadrant are closed orbits: both populations oscillate periodically, with predator oscillations lagging behind prey oscillations.

**Linearization at $(0,0)$.** The Jacobian at the origin is $J(0,0) = \begin{pmatrix} 1 & 0 \\ 0 & -1 \end{pmatrix}$ with eigenvalues $1$ and $-1$: a \textbf{saddle}. The $x$-axis (prey only, no predators) and $y$-axis (predators only, no prey) are invariant; along the $x$-axis prey grow exponentially, along the $y$-axis predators die off exponentially.


\bigskip


\section{Nonlinear Systems and Linearization}\label{sec:3.5}

\subsection{The Jacobian at a Fixed Point}\label{sec:3.5.1}

Consider a general system $\dot{x} = f(x)$ with $f: \mathbb{R}^n \to \mathbb{R}^n$ of class $C^1$, and suppose $x^*$ is a fixed point: $f(x^*) = 0$. Writing $u = x - x^*$ for the displacement from the fixed point, Taylor expansion gives

\[
\dot{u} = f(x^* + u) = \underbrace{f(x^*)}_{=0} + Df(x^*)\, u + O(\|u\|^2),
\]

where $Df(x^*)$ is the \textbf{Jacobian matrix} of $f$ at $x^*$:

\[
Df(x^*)_{ij} = \frac{\partial f_i}{\partial x_j}\bigg|_{x = x^*}.
\]

The \textbf{linearization} of the system at $x^*$ is the linear system $\dot{u} = Df(x^*)\, u$. The natural question is: to what extent does the linearization faithfully represent the local dynamics of the nonlinear system?

\subsection{Hyperbolic Fixed Points and the Hartman--Grobman Theorem}\label{sec:3.5.2}

\begin{definition}\label{def:3.7.}
A fixed point $x^*$ is called \textbf{hyperbolic} if all eigenvalues of $Df(x^*)$ have nonzero real part.

Equivalently, no eigenvalue lies on the imaginary axis. This excludes centers ($\operatorname{Re}(\lambda) = 0$) and cases with a zero eigenvalue.
\end{definition}


\textbf{Theorem 3.8 (Hartman--Grobman).} *Let $x^*$ be a hyperbolic fixed point of $\dot{x} = f(x)$ where $f$ is $C^1$. Then there exists a neighborhood $U$ of $x^*$ and a homeomorphism $h: U \to V$ (where $V$ is a neighborhood of the origin) that maps trajectories of the nonlinear system to trajectories of the linearized system $\dot{u} = Df(x^*)\,u$, preserving the direction of time.*

In precise terms, $h$ conjugates the flow: $h(\phi_t(x)) = e^{Df(x^*)t} h(x)$ for $x$ near $x^*$ and small $|t|$.

\textbf{Significance.} The theorem says that near a hyperbolic fixed point, the nonlinear flow is \textit{topologically equivalent} to its linearization. The phase portrait of the nonlinear system looks qualitatively the same as that of the linear system. In particular:

\begin{itemize}[nosep]
  \item A stable node of the linearization is a stable node of the nonlinear system.
  \item A saddle of the linearization is a saddle of the nonlinear system.
  \item A stable spiral of the linearization is a stable spiral of the nonlinear system.
  \item And so on.
\end{itemize}

\textbf{Limitations.} The theorem has two important limitations:

\begin{enumerate}[nosep]
  \item **The conjugacy $h$ is only a homeomorphism, not a diffeomorphism.** It preserves the topology of the phase portrait but may distort angles, smoothness of trajectories, etc. (In fact, a smooth version---the Sternberg linearization theorem---holds under additional non-resonance conditions on the eigenvalues, but the general statement requires only continuity.)

  \item \textbf{The theorem says nothing about non-hyperbolic fixed points.} When the Jacobian has eigenvalues on the imaginary axis, the nonlinear terms can qualitatively change the dynamics. The most important example is the \textbf{center}: the linearization of the simple pendulum at $(0,0)$ is a center, and it happens that the nonlinear system also has a center there (due to the conserved energy), but this is \textit{not} guaranteed in general. A small perturbation can turn a linear center into a spiral. The determination of stability for non-hyperbolic fixed points requires more refined tools, such as center manifold theory or Lyapunov functions.
\end{enumerate}

\subsection{Stable and Unstable Manifolds}\label{sec:3.5.3}

\textbf{Theorem 3.9 (Stable Manifold Theorem).} *Let $x^*$ be a hyperbolic fixed point of $\dot{x} = f(x)$ with $f \in C^1$. Let $E^s$ and $E^u$ denote the stable and unstable subspaces of the linearization (the eigenspaces corresponding to eigenvalues with negative and positive real parts, respectively). Then there exist $C^1$ embedded submanifolds $W^s(x^*)$ and $W^u(x^*)$ of $\mathbb{R}^n$ such that:*

\begin{enumerate}[nosep]
  \item *$W^s(x^*)$ is tangent to $E^s$ at $x^*$, and $\dim W^s = \dim E^s$.*
  \item *$W^u(x^*)$ is tangent to $E^u$ at $x^*$, and $\dim W^u = \dim E^u$.*
  \item *$W^s(x^*) = \{x : \phi_t(x) \to x^* \text{ as } t \to +\infty\}$ (locally, near $x^*$).*
  \item *$W^u(x^*) = \{x : \phi_t(x) \to x^* \text{ as } t \to -\infty\}$ (locally, near $x^*$).*
  \item *$W^s$ and $W^u$ are invariant under the flow.*
\end{enumerate}

In the linear case ($f(x) = Ax$), the stable and unstable manifolds are exactly $E^s$ and $E^u$ (linear subspaces). In the nonlinear case, they are curved, but they are tangent to the corresponding linear subspaces at the fixed point.

\textbf{Worked Example 3.10.} Consider the system

\[
\dot{x}_1 = -x_1, \qquad \dot{x}_2 = x_2 + x_1^2.
\]

The origin is a fixed point. The Jacobian is $Df(0) = \begin{pmatrix} -1 & 0 \\ 0 & 1 \end{pmatrix}$, a saddle with $E^s = \{x_2 = 0\}$ and $E^u = \{x_1 = 0\}$.

The stable manifold $W^s$: on it, $\phi_t(x) \to 0$ as $t \to \infty$. From the first equation, $x_1(t) = c_1 e^{-t}$. Substituting into the second: $\dot{x}_2 = x_2 + c_1^2 e^{-2t}$. This linear ODE has solution $x_2(t) = c_2 e^t - \frac{c_1^2}{3}e^{-2t}$. For $x_2(t) \to 0$ as $t \to \infty$, we need $c_2 = 0$, giving $x_2(t) = -\frac{c_1^2}{3}e^{-2t}$. Since $x_1(t) = c_1 e^{-t}$, we have $x_1^2 = c_1^2 e^{-2t}$ and therefore $x_2 = -x_1^2/3$.

The local stable manifold is the parabola $W^s_{\text{loc}} = \{(x_1, x_2) : x_2 = -x_1^2/3\}$, which is tangent to $E^s = \{x_2 = 0\}$ at the origin, as the theorem predicts.

Similarly, one can compute that the unstable manifold is exactly the $x_2$-axis: $W^u = \{x_1 = 0\}$ (since $x_1 = 0$ is invariant and $\dot{x}_2 = x_2$ there).


\bigskip


\section{Limit Cycles and the Poincar\'e--Bendixson Theorem}\label{sec:3.6}

\subsection{Limit Cycles}\label{sec:3.6.1}

A \textbf{limit cycle} is an isolated periodic orbit. "Isolated" means that nearby trajectories are not periodic---they spiral toward or away from the limit cycle.

\begin{itemize}[nosep]
  \item A \textbf{stable limit cycle} (attracting) is one where nearby trajectories spiral toward it as $t \to +\infty$.
  \item An \textbf{unstable limit cycle} (repelling) is one where nearby trajectories spiral away.
  \item A \textbf{semi-stable limit cycle} attracts from one side and repels from the other.
\end{itemize}

Limit cycles are inherently nonlinear phenomena: linear systems in 2D can have periodic orbits (centers), but these are never isolated---they come in continuous families filling an annular region.

\subsection{The Poincar\'e--Bendixson Theorem}\label{sec:3.6.2}

The following theorem is one of the most important results in planar dynamics. It places severe constraints on what the long-term behavior of a two-dimensional flow can look like.

\textbf{Theorem 3.11 (Poincar\'e--Bendixson).} *Let $\dot{x} = f(x)$ be a $C^1$ system on an open subset of $\mathbb{R}^2$. Suppose $\gamma^+$ is a forward trajectory that remains in a compact set $K$ containing only finitely many fixed points. Then the $\omega$-limit set $\omega(\gamma^+)$ is one of the following:*

\begin{enumerate}[nosep]
  \item \textit{A single fixed point.}
  \item \textit{A periodic orbit.}
  \item \textit{A finite union of fixed points together with trajectories connecting them (a heteroclinic or homoclinic cycle).}
\end{enumerate}

\begin{proof}[Proof sketch]
sketch.** The key geometric insight is that in two dimensions, a trajectory cannot cross itself (by uniqueness of solutions) and hence must "use up" the available space in a monotone way. More precisely, consider a local cross-section (a short line segment transverse to the flow). The return map to this section (if the trajectory returns) must be monotone. Monotone maps on intervals have simple dynamics: every orbit converges. This forces the $\omega$-limit set to be one of the three types listed. The rigorous proof uses the Jordan curve theorem and requires some careful topological arguments; see Hirsch, Smale, and Devaney (2013), Chapter 10, or Perko (2001), Chapter 3.

\textbf{Corollary 3.12.} *If $K$ is a compact, positively invariant region containing no fixed points, then $K$ contains a periodic orbit.*

\begin{proof}
* Take any trajectory in $K$. It remains in $K$ for all forward time (positive invariance), and $K$ is compact, so the trajectory remains in a compact set. By Poincar\'e--Bendixson, the $\omega$-limit set is a fixed point, a periodic orbit, or a heteroclinic cycle. Since $K$ has no fixed points, options (1) and (3) are excluded. $\blacksquare$

This corollary provides a powerful technique for \textit{proving the existence} of limit cycles: construct a trapping region with no fixed points.

\end{proof}

\end{proof}

\subsection{No Chaos in Two Dimensions}\label{sec:3.6.3}

The Poincar\'e--Bendixson theorem has a remarkable consequence:

\textbf{Corollary 3.13.} *A continuous-time autonomous system on $\mathbb{R}^2$ (or any subset of $\mathbb{R}^2$) cannot exhibit chaotic behavior.*

The argument is straightforward: chaotic dynamics requires trajectories to have complicated $\omega$-limit sets (e.g., strange attractors), but Poincar\'e--Bendixson forces $\omega$-limit sets in 2D to be one of three simple types. In particular, sensitive dependence on initial conditions in the sense of chaos is impossible.

This is why one needs at least three dimensions for continuous-time chaos. The Lorenz system, R\"ossler attractor, and other classic chaotic flows all live in $\mathbb{R}^3$ or higher. This dimensional threshold is a fundamental difference between continuous and discrete dynamics, where even one-dimensional maps (like the logistic map from Chapter 2) can exhibit chaos.

\subsection{The Van der Pol Oscillator}\label{sec:3.6.4}

The \textbf{Van der Pol oscillator} is a second-order ODE arising from electrical circuit theory:

\begin{equation}\label{eq:3.9}\tag{3.9}
\ddot{x} + \mu(x^2 - 1)\dot{x} + x = 0, \quad \mu > 0.
\end{equation}

Note the damping term $\mu(x^2 - 1)\dot{x}$: when $|x| < 1$, the coefficient $(x^2 - 1)$ is negative, so the "damping" is actually \textit{negative} (energy is pumped in); when $|x| > 1$, the damping is positive (energy is dissipated). This balance gives rise to a stable limit cycle.

Converting to a first-order system with $x_1 = x$, $x_2 = \dot{x}$:

\[
\dot{x}_1 = x_2, \qquad \dot{x}_2 = -x_1 + \mu(1 - x_1^2) x_2.
\]

\textbf{Fixed point.} The only fixed point is the origin. The Jacobian there is

\[
Df(0) = \begin{pmatrix} 0 & 1 \\ -1 & \mu \end{pmatrix}.
\]

The eigenvalues satisfy $\lambda^2 - \mu\lambda + 1 = 0$, so $\lambda = \frac{\mu \pm \sqrt{\mu^2 - 4}}{2}$. For $0 < \mu < 2$, the eigenvalues are complex with positive real part $\mu/2 > 0$: the origin is an \textbf{unstable spiral}. For $\mu \geq 2$, both eigenvalues are real and positive: an \textbf{unstable node}.

In either case, trajectories are repelled from the origin. One can show (and we take this on faith, though the argument uses a trapping region and the Poincar\'e--Bendixson theorem) that trajectories starting far from the origin are pushed inward: the dissipative damping for $|x| > 1$ eventually dominates. By constructing a compact annular region that trajectories cannot leave, and noting that it contains no fixed points, Corollary 3.12 guarantees the existence of a periodic orbit.

For small $\mu$, the limit cycle is approximately circular with radius 2 (this can be shown by the method of averaging). For large $\mu$, the limit cycle develops sharp transitions and slow-fast structure, approaching a \textit{relaxation oscillation}.


\bigskip


\section{Lyapunov Functions}\label{sec:3.7}

\subsection{Lyapunov Stability}\label{sec:3.7.1}

We have seen several methods for determining the stability of a fixed point: eigenvalue analysis for linear systems, Hartman--Grobman for hyperbolic fixed points of nonlinear systems. But what about non-hyperbolic fixed points? And what about \textit{global} (as opposed to local) stability? Lyapunov functions provide a powerful and flexible approach.

\begin{definition}\label{def:3.14.}
Let $x^*$ be a fixed point of $\dot{x} = f(x)$. A $C^1$ function $V: U \to \mathbb{R}$ defined on a neighborhood $U$ of $x^*$ is called a \textbf{Lyapunov function} if:

1. $V(x^*) = 0$ and $V(x) > 0$ for all $x \in U \setminus \{x^*\}$.
2. $\dot{V}(x) \leq 0$ for all $x \in U$, where $\dot{V}$ denotes the time derivative along trajectories:

\[
\dot{V}(x) = \nabla V(x) \cdot f(x) = \sum_{i=1}^n \frac{\partial V}{\partial x_i} f_i(x).
\]

Note that $\dot{V}$ can be computed without solving the ODE; it depends only on $V$ and $f$.
\end{definition}


\textbf{Theorem 3.15 (Lyapunov Stability Theorem).} *Let $x^*$ be a fixed point of $\dot{x} = f(x)$ and let $V$ be a Lyapunov function on a neighborhood of $x^*$.*

\begin{enumerate}[nosep]
  \item *If $\dot{V}(x) \leq 0$ in $U$, then $x^*$ is \textbf{Lyapunov stable} (trajectories starting near $x^*$ remain near $x^*$).*
  \item *If $\dot{V}(x) < 0$ for all $x \in U \setminus \{x^*\}$, then $x^*$ is \textbf{asymptotically stable} (trajectories starting near $x^*$ converge to $x^*$).*
\end{enumerate}

\begin{proof}[Proof of (1)]
of (1).* Let $\varepsilon > 0$ be small enough that $\overline{B}_\varepsilon(x^*) \subset U$. Since $V$ is continuous, positive on the compact set $S_\varepsilon = \{x : \|x - x^*\| = \varepsilon\}$, and $S_\varepsilon$ does not contain $x^*$, we have $m = \min_{S_\varepsilon} V > 0$. By continuity of $V$ and $V(x^*) = 0$, there exists $\delta > 0$ such that $V(x) < m$ for all $x \in B_\delta(x^*)$.

Now take any trajectory starting in $B_\delta(x^*)$. Since $\dot{V} \leq 0$, $V$ is non-increasing along the trajectory, so $V(x(t)) \leq V(x(0)) < m$ for all $t \geq 0$. But $V \geq m$ on $S_\varepsilon$, so the trajectory can never reach $S_\varepsilon$. Therefore $x(t) \in B_\varepsilon(x^*)$ for all $t \geq 0$.
\end{proof}


\begin{proof}[Proof of (2)]
of (2).* By (1), $x^*$ is Lyapunov stable. It remains to show convergence. Given the stronger hypothesis $\dot{V} < 0$, the function $V(x(t))$ is strictly decreasing, hence converges to some limit $\ell \geq 0$. If $\ell > 0$, then the trajectory stays in the compact annular region $\{x : \ell \leq V(x) \leq V(x(0))\}$, which does not contain $x^*$. On this compact set, $\dot{V}$ achieves a maximum $-c < 0$ (since $\dot{V}$ is continuous and strictly negative away from $x^*$). Then $V(x(t)) \leq V(x(0)) - ct \to -\infty$, contradicting $V \geq 0$. So $\ell = 0$ and $V(x(t)) \to 0$, which forces $x(t) \to x^*$. $\blacksquare$

\end{proof}

\subsection{Example: Energy as a Lyapunov Function}\label{sec:3.7.2}

Consider the damped pendulum

\[
\ddot{\theta} + b\dot{\theta} + \sin\theta = 0, \quad b > 0,
\]

or equivalently $\dot{x}_1 = x_2$, $\dot{x}_2 = -\sin x_1 - b x_2$, with the equilibrium at $(0,0)$.

Define the energy $V(x_1, x_2) = \frac{1}{2}x_2^2 + (1 - \cos x_1)$.

\begin{itemize}[nosep]
  \item $V(0,0) = 0$, and $V(x_1, x_2) > 0$ for $(x_1, x_2)$ in a neighborhood of the origin (since $1 - \cos x_1 \geq 0$ with equality only at $x_1 = 2k\pi$, and $\frac{1}{2}x_2^2 \geq 0$; both vanish simultaneously only at the origin, within $|x_1| < 2\pi$).

  \item $\dot{V} = x_2 \dot{x}_2 + \sin(x_1)\dot{x}_1 = x_2(-\sin x_1 - bx_2) + \sin(x_1) \cdot x_2 = -bx_2^2 \leq 0$.
\end{itemize}

So $V$ is a Lyapunov function, and the origin is Lyapunov stable.

Can we conclude asymptotic stability? We have $\dot{V} = -bx_2^2$, which is zero not only at the origin but on the entire $x_1$-axis (where $x_2 = 0$). So condition (2) of the theorem is not met. However, a trajectory cannot remain on $\{x_2 = 0\}$ unless it is at a fixed point, because if $x_2 = 0$ and $x_1 \neq 0$ (and $x_1 \neq k\pi$), then $\dot{x}_2 = -\sin x_1 \neq 0$ forces the trajectory off the $x_1$-axis. This is the idea behind \textbf{LaSalle's invariance principle}, which strengthens the Lyapunov theorem:

\textbf{Theorem 3.16 (LaSalle's Invariance Principle).} *Under the hypotheses of Theorem 3.15(1), let $S = \{x \in U : \dot{V}(x) = 0\}$, and let $M$ be the largest invariant subset of $S$. Then every trajectory starting in a compact level set $\{V \leq c\} \subset U$ converges to $M$ as $t \to \infty$.*

For the damped pendulum, $S = \{x_2 = 0\}$ and the only invariant subset of $S$ near the origin is $M = \{(0,0)\}$. By LaSalle's principle, $(0,0)$ is asymptotically stable.

\subsection{Remarks on Finding Lyapunov Functions}\label{sec:3.7.3}

There is no general algorithm for constructing Lyapunov functions, and this is the main practical difficulty. However, physical intuition often helps:

\begin{itemize}[nosep]
  \item \textbf{Mechanical systems:} Total energy (kinetic + potential) is a natural candidate.
  \item \textbf{Gradient systems} $\dot{x} = -\nabla V(x)$: the function $V$ itself is a Lyapunov function, since $\dot{V} = \nabla V \cdot (-\nabla V) = -\|\nabla V\|^2 \leq 0$.
  \item \textbf{Quadratic forms:} For linear systems $\dot{x} = Ax$ with $A$ stable (all eigenvalues in the left half-plane), the Lyapunov equation $A^T P + PA = -Q$ has a unique positive-definite solution $P$ for any positive-definite $Q$, and $V(x) = x^T P x$ is a Lyapunov function.


\bigskip

\end{itemize}

\section{Looking Ahead}\label{sec:3.8}

In this chapter we have developed the theory of continuous dynamical systems on $\mathbb{R}^n$, with particular attention to the plane. We saw that the Poincar\'e--Bendixson theorem constrains planar flows to have simple $\omega$-limit sets, ruling out chaos. In Chapter 4, we will venture into three and higher dimensions, where the Poincar\'e--Bendixson theorem no longer applies and the door opens to chaotic attractors, sensitive dependence on initial conditions, and the rich phenomena that motivate ergodic theory.


\bigskip


\section*{Exercises}

\begin{exercise}
Consider the linear system $\dot{x} = Ax$ where $A = \begin{pmatrix} 1 & -5 \\ 1 & -3 \end{pmatrix}$.

(a) Find the eigenvalues and classify the fixed point at the origin.

(b) Find the general solution using the matrix exponential.

(c) Sketch the phase portrait.
\end{exercise}


\begin{exercise}
(Failure of uniqueness.) Show that the initial value problem $\dot{x} = \sqrt{|x|}$, $x(0) = 0$ has infinitely many solutions. Verify that each of the functions

\[
x_a(t) = \begin{cases} 0 & t \leq a, \\ \frac{1}{4}(t - a)^2 & t > a \end{cases}
\]

is a solution for every $a \geq 0$. Why does the Picard--Lindel\"of theorem not apply?
\end{exercise}


\begin{exercise}
(Nullcline analysis.) For the system

\[
\dot{x} = x(3 - x - 2y), \qquad \dot{y} = y(2 - x - y),
\]

(a) Find all fixed points in the first quadrant (including the boundary).

(b) Sketch the nullclines and determine the direction of the flow in each region.

(c) Determine the stability of each fixed point using linearization.

(d) Sketch the phase portrait in the first quadrant.
\end{exercise}


\begin{exercise}
(Lyapunov function.) Consider the system $\dot{x}_1 = -x_1 + x_1 x_2^2$, $\dot{x}_2 = -x_2 + x_1^2 x_2$. Show that $V(x_1, x_2) = x_1^2 + x_2^2$ is a Lyapunov function in the region $\{x_1^2 + x_2^2 < 1\}$ and conclude that the origin is asymptotically stable.
\end{exercise}


\begin{exercise}
(Stable manifold computation.) For the system

\[
\dot{x}_1 = x_1 + x_2^2, \qquad \dot{x}_2 = -x_2,
\]

(a) Find the linearization at the origin and identify $E^s$ and $E^u$.

(b) Compute the stable manifold $W^s$ exactly (proceed as in Worked Example 3.10).

(c) Compute the unstable manifold $W^u$ by a power-series expansion $x_2 = h(x_1) = a x_1 + b x_1^2 + \cdots$ and matching coefficients.
\end{exercise}


\begin{exercise}
(Poincar\'e--Bendixson in action.) Consider the system in polar coordinates

\[
\dot{r} = r(1-r^2) + \mu r \cos\theta, \qquad \dot{\theta} = 1,
\]

for small $\mu > 0$.

(a) Show that for $\mu = 0$, the unit circle $r = 1$ is a stable limit cycle.

(b) For small $\mu > 0$, construct a trapping region (an annulus $r_1 \leq r \leq r_2$ with appropriate $r_1, r_2$) and use the Poincar\'e--Bendixson theorem to prove that a periodic orbit still exists.
\end{exercise}


\begin{exercise}
(Jordan blocks and polynomial growth.) Consider the $3 \times 3$ system $\dot{x} = Jx$ where

\[
J = \begin{pmatrix} 0 & 1 & 0 \\ 0 & 0 & 1 \\ 0 & 0 & 0 \end{pmatrix}.
\]

(a) Compute $e^{Jt}$ directly from the series definition (note $J$ is nilpotent: find the smallest $k$ with $J^k = 0$).

(b) Write down the general solution and describe the trajectories. Show that $\|x(t)\|$ grows polynomially (not exponentially) in time.

(c) Is the origin stable? Lyapunov stable? Asymptotically stable?
\end{exercise}


\begin{exercise}
(Lotka--Volterra conserved quantity.) For the system (3.8), verify that $H(x,y) = x - \ln x + y - \ln y$ satisfies $\dot{H} = 0$ along trajectories. Conclude that all orbits in the open first quadrant are periodic.

---
\end{exercise}


\section*{Recommended Reading}

\begin{itemize}[nosep]
  \item \textbf{Hirsch, M. W., Smale, S., and Devaney, R. L.} \textit{Differential Equations, Dynamical Systems, and an Introduction to Chaos}, 3rd edition. Academic Press, 2013. An excellent undergraduate text; Chapters 1--10 cover all the material in this chapter and more, with many examples and clear exposition.

  \item \textbf{Strogatz, S. H.} \textit{Nonlinear Dynamics and Chaos}, 2nd edition. Westview Press, 2015. The standard "first course" in nonlinear dynamics. Highly readable, with wonderful physical intuition and examples. Chapters 5--8 are particularly relevant.

  \item \textbf{Perko, L.} \textit{Differential Equations and Dynamical Systems}, 3rd edition. Springer, 2001. More rigorous and comprehensive than Strogatz; includes proofs of the Hartman--Grobman theorem, the stable manifold theorem, and the Poincar\'e--Bendixson theorem. A good reference for the reader wanting complete proofs.

  \item \textbf{Arnold, V. I.} \textit{Ordinary Differential Equations}, 3rd edition. Springer, 1992. Arnold's geometric approach to ODEs is unique and illuminating. The treatment of flows and the matrix exponential in Chapters 1--5 is particularly elegant. A challenging but rewarding read.

  \item \textbf{Coddington, E. A., and Levinson, N.} \textit{Theory of Ordinary Differential Equations}. McGraw-Hill, 1955. The classic graduate reference for existence, uniqueness, and continuation theorems. Included here for the reader wishing to see the proof of Peano's theorem and other foundational results.
\end{itemize}
